Consideramos por último, el caso de 4 $e^-$ equivalentes con $l=2$. La correspondiente configuración electrónica es $d^{4}$. En este caso $\mathcal{D}\in [0,2]$ y las reglas son \\
\underline{\textbf{Reglas}}:
\begin{align}
  &\mathcal{D}=2 \longrightarrow 1\times (M_{S}=0):\ (+-+-)\\
  &\mathcal{D}=1 \longrightarrow \begin{cases}
    1\times (M_{S}=1):\ (+-++)\\
    2\times (M_{S}=0):\ (+-+-)
  \end{cases}\\
  &\mathcal{D}=0 \longrightarrow \begin{cases}
    1\times (M_{S}) = 2:\ (++++)\\
    4\times (M_{S}=1):\ (+++-)\\
    6\times (M_{S}=0):\ (++--)
  \end{cases}
 \end{align}
Entonces los términos espectrales son 


\begin{align*}
 \begin{array}{c | c | c | c}
   M_{L} & \text{ Particiones } & \text{ Duplas } & \text{ Términos }\\
   \hline
   M_{L}= 6  & 6 = 2 + 2 + 1 + 1  &  \mathcal{D} = 2 & 1 \times (M_{S} = 0) \implies \ce{^{1}I}(13)\\
   \hline
   M_{L} = 5 & 5 = 2 + 2 + 1 + 0 & \mathcal{D} = 1 & \begin{dcases} 1 \times (M_{S} = 1) \implies \ce{^{3}H}(33)\\
     2\times (M_{S}=0)\implies \cancel{\ce{^{1}I}(13)},\ \cancel{\ce{^{3}H(33)}} 
   \end{dcases}\\
  \hline
     M_{L} = 4 & 
     \begin{aligned}
       4 &= 2 + 2  + 0 + 0\\ 
         &= 2 + 1 + 1 + 0\\
         &= 2 + 2 + 1 - 1
     \end{aligned} 
     & \begin{aligned}
       \mathcal{D} &= 2\\ 
       \mathcal{D} &= 1\\
       \mathcal{D} &=1
     \end{aligned} & \begin{dcases}
      2 \times (M_{S} = 1) \\
      5\times (M_{S} = 0) 
   \end{dcases}  \implies \ce{^{3}G(27)},\ 2\times\ce{^{1}G(9)} \\ 
     \hline
     M_{L} = 3 & 
     \begin{aligned}
     3 &= 2 + 1 + 0 + 0\\ 
       &= 2 + 1 + 1 - 1\\
       &= 2 + 2 - 1 + 0\\
       &= 2 + 2 + 1 -2 
   \end{aligned} & 
   \begin{aligned}
      \mathcal{D} &= 1\\ 
      \mathcal{D} &= 1\\
      \mathcal{D}&= 1\\
      \mathcal{D}&=1 
   \end{aligned} & 
    \begin{cases}
    4 \times (M_{S} = 1)\\
    8 \times (M_{S} = 0) 
  \end{cases} \implies  2\times\ce{^{3}F}(21),\ \ce{^{1}F}(7)\\
 \hline
   M_{L} = 2 & 
     \begin{aligned}
       2 &= 1 + 1 + 0 + 0\\ 
         &= 2 + 1 - 1 + 0\\
         &= 2 + 1 + 1 - 2\\
         &= 2 + 2 + 0 - 2\\
         &= 2 + 2 - 1 - 1
     \end{aligned} 
     & \begin{aligned}
       \mathcal{D} &= 2\\ 
       \mathcal{D} &= 0\\
       \mathcal{D} &= 1\\
       \mathcal{D} &= 1\\
       \mathcal{D} &= 2
     \end{aligned} & \begin{dcases}
      1 \times (M_{S} = 2) \\
      6\times(M_{S}=1)\\ 
      12\times (M_{S} = 0) 
   \end{dcases}  \implies \ce{^{5}D(25)}, \ce{^{3}D}(15),\;  2\times \ce{^{1}D(5)} \\ 
     \hline
     M_{L} = 1 & 
     \begin{aligned}
     1 &= 1 + 1 - 1 + 0\\
       &= 2 + 0 + 0 - 1\\
       &= 2 + 1 - 1 - 1\\
       &= 2 + 1 + 0 - 2\\
       &= 2 + 2 - 1 - 2
   \end{aligned} & 
   \begin{aligned}
      \mathcal{D} &= 1\\
      \mathcal{D} &= 1\\
      \mathcal{D} &= 1\\
      \mathcal{D} &= 0\\
      \mathcal{D} &= 1 
   \end{aligned} & 
    \begin{cases}
    1\times (M_{S}=2)\\  
    8 \times (M_{S} = 1)\\
    14 \times (M_{S} = 0) 
    \end{cases} \implies  2\times\ce{^{3}P}(9),\ \\
 \hline 
   M_{L} = 0 & 
     \begin{aligned}
       0 &= 1 + 1 - 1 - 1\\
         &= 1 + 1 + 0 - 2\\
         &= 1 + 0 + 0 - 1\\
         &= 2 + 0 - 1 - 1\\
         &= 2 + 0 + 0 - 2\\
         &= 2 + 1 - 1 - 2\\
         &= 2 + 2 - 2 - 2 
     \end{aligned} 
     & \begin{aligned}
       \mathcal{D} &= 2\\ 
       \mathcal{D} &= 1\\
       \mathcal{D} &= 1\\
       \mathcal{D} &= 1\\
       \mathcal{D} &= 1\\
       \mathcal{D} &= 0\\
       \mathcal{D} &= 2
     \end{aligned} & \begin{dcases}
      1 \times (M_{S} = 2) \\
      8\times (M_{S} = 1)\\
      16\times(M_{S}=0) 
   \end{dcases}  \implies 2\times\ce{^{1}S(1)} \\ 
      \end{array}
\end{align*}
\begin{align*}
 g_{T} = 13 + 33 + 29 + 2 \times 9 + 2 \times 21 + 7 + 25 + 15 + 5 + 2 \times 9 + 2 = 210 = \binom{10}{4}
\end{align*}
Por tanto, los términos L-S son: 
\begin{align*}
  \boxed{\text{ T.E.: } \begin{aligned}[t]
    &\ce{^{1}I}(13),\;  \ce{^{3}H}(33),\;  \ce{^{3}G}(27),\;  2\times\ce{^{1}G}(9),\;  2\times\ce{^{3}F}(21),\;  \ce{^{1}F}(7),\;  \ce{^{5}D}(25),\;  \ce{^{3}D}(15),\\ 
    &\quad 2\times\ce{^{1}D}(5),\;  2\times\ce{^{3}P}(9),\;  2\times\ce{^{1}S}(1)
  \end{aligned}}
\end{align*}
\textbf{Nota:} Operativamente, es conveniente asignar un vector a cada regla de la siguiente forma: 
\begin{align*}
  \mathcal{D} \to  (\mathrm{mult.}( M_{S} = 2),\; \mathrm{mult.} (M_{S} = 1),\; \mathrm{mult.} (M_{S} = 0)) \implies 
  \begin{dcases}
    \mathcal{D} = 2 \to  (0,0,1)\\ 
    \mathcal{D} = 1 \to  (0,1,2)\\ 
    \mathcal{D} = 0 \to  (1,4,6)
  \end{dcases}
\end{align*}
De esta forma, para un cierto $M_{L}$ se le asigna un vector con el que podemos operar para obtener los términos espectrales. Por ejemplo, para $M_{L} = 2$ tendremos el vector asociado:
\begin{align*}
2 \times (0,0,1) + 2 \times (0,1,2) + 1 \times (1,4,6) = (1,6,12)
\end{align*}
donde los coeficientes de los vectores vienen dados por el número de veces que una regla se cumple. Así, podemos restar a este vector el asociado al anterior $M_{L}$, es decir, $M_{L} = 3$ y obtenemos: 
\begin{align*}
(1,6,12) - (0,4,8) = (1,2,4)
\end{align*}
Y este vector es lo que necesitamos para obtener los términos espectrales emergentes. Con la primera componente podemos asegurar que emerge un término con $S = 2$. De la segunda componente (2), al restar su anterior componente (1), nos da un término emergente con $S = 1$. Y con la tercera componente (4), al restar su anterior componente (2) nos da 2 términos emergentes con $S = 0$. Este mecanismo de resta incorpora las consideraciones de que términos con $S$ superior también han de tenerse en cuenta en la contabilidad de $M_{S}$ inferiores.   
